\section{Conclusion}

\noindent This paper presented the IMU-Controlled UltraWideband Mesh (I.C.U.M.), a decentralized localization framework designed to enable autonomous swarm compliance in GPS-denied environments. By integrating low-cost inertial sensors with peer-to-peer Ultra-Wideband ranging, we demonstrated that mobile agents can establish and maintain a rigid relative coordinate frame without reliance on static anchors or absolute positioning systems.

\noindent Our empirical results confirm that the proposed Event-Triggered Messaging (ETM) protocol successfully balances the trade-off between formation integrity and energy efficiency. By utilizing IMU data to gate radio transmissions, the system achieved a 76.17\% reduction in communication overhead compared to periodic baselines, while maintaining sub-meter relative accuracy (RMSE $\approx$ 0.1m in high-density clusters). Furthermore, we identified the critical role of node density in formation stability, observing a "Rigidity Transition" at $N \approx 35$, below which the network remains geometrically unstable.

\noindent These findings have significant implications for the deployment of search-and-rescue swarms and subterranean robotics. We have shown that by prioritizing relative geometric rigidity over absolute latitude, autonomous teams can operate safely and cohesively in completely unknown environments. While scalability limitations regarding channel saturation ($O(N^2)$) remain a challenge for future work, the I.C.U.M. architecture provides a verified foundation for the next generation of resilient, energy-aware distributed systems.
